\section*{Hoofdstuk \ref{ch:g11:functional-groups} --- Karakteristieke groepen en families}

\begin{solution}{exo:g11:functional-groups:1}
Hydroxyl, \omterm{def:g11:functional-groups:alcohol}{alcohol}; carbonyl midden in de keten, \omterm{def:g11:functional-groups:carbonyl}{keton}; carboxyl,
\omterm{def:g11:functional-groups:carboxylic-acid}{carbonzuur}; estergroep, \omterm{def:g11:functional-groups:carboxylic-acid}{ester}; aminogroep \ce{-NH2}, \omterm{def:g11:functional-groups:amine}{amine}.
\end{solution}

\begin{solution}{exo:g11:functional-groups:2}
Methanol; propaan-1-ol; methaanzuur; propanal; propanon.
\end{solution}

\begin{solution}{exo:g11:functional-groups:3}
\omterm{def:g11:functional-groups:alcohol}{Alcohol}; \omterm{def:g11:functional-groups:carbonyl}{aldehyde}; \omterm{def:g11:functional-groups:carboxylic-acid}{ester}; \omterm{def:g11:functional-groups:amine}{amide}; \omterm{def:g11:functional-groups:halogenoalkane}{halogeenalkaan}; \omterm{def:g11:functional-groups:halogenoalkane}{alkeen}.
\end{solution}

\begin{solution}{exo:g11:functional-groups:4}
Propanal \ce{CH3-CH2-CHO}, propanon \ce{CH3-CO-CH3}. Propanal is het
\omterm{def:g11:functional-groups:carbonyl}{aldehyde}: het \ce{C=O}-koolstofatoom zit aan het uiteinde van de keten,
gebonden aan een waterstofatoom; in propanon zit het tussen twee
koolstofatomen.
\end{solution}

\begin{solution}{exo:g11:functional-groups:5}
Primair (het koolstofatoom met de \ce{-OH} heeft één koolstofbuur);
secundair (twee); tertiair (drie).
\end{solution}

\begin{solution}{exo:g11:functional-groups:6}
\setchemfig{atom sep=1.5em}
Pentaan-3-ol \chemfig{-[:30]-[:-30](-[:-90]OH)-[:30]-[:-30]};
2-methylbutaanzuur \chemfig{-[:30]-[:-30](-[:-90])-[:30](=[:90]O)-[:-30]OH};
hexaan-2-on \chemfig{-[:30](=[:90]O)-[:-30]-[:30]-[:-30]-[:30]};
propylethanoaat \chemfig{-[:30](=[:90]O)-[:-30]O-[:30]-[:-30]-[:30]}.
\end{solution}

\begin{solution}{exo:g11:functional-groups:7}
\ce{C3H6O}. Ja, het zijn \omterm{def:g11:organic-skeletons:isomer}{isomeren}, maar met verschillende groepen: een
\omterm{def:g11:functional-groups:carbonyl}{aldehyde} en een \omterm{def:g11:functional-groups:carbonyl}{keton} (allebei carbonylgroepen, op verschillende
plaatsen).
\end{solution}

\begin{solution}{exo:g11:functional-groups:8}
\ce{HCOO-CH2-CH3}, ethylmethanoaat; \ce{CH3-COO-CH2-CH2-CH3},
propylethanoaat.
\end{solution}

\begin{solution}{exo:g11:functional-groups:9}
Methylethanoaat, uit ethaanzuur en methanol; ethylpropanoaat, uit
propaanzuur en ethanol; methylmethanoaat, uit methaanzuur en methanol.
\end{solution}

\begin{solution}{exo:g11:functional-groups:10}
\setchemfig{atom sep=1.5em}
\begin{center}
\small
\begin{tabular}{cccc}
\chemfig{-[:30]-[:-30]-[:30]-[:-30]OH} &
\chemfig{-[:30](-[:90]OH)-[:-30]-[:30]} &
\chemfig{-[:30](-[:90])-[:-30]-[:30]OH} &
\chemfig{-[:0](-[:90]OH)(-[:-90])-[:0]} \\
butaan-1-ol & butaan-2-ol & 2-methylpropaan-1-ol & 2-methylpropaan-2-ol \\
primair & secundair & primair & tertiair
\end{tabular}
\end{center}
\end{solution}

\begin{solution}{exo:g11:functional-groups:11}
Een \omterm{def:g11:functional-groups:amine}{amide} (ethaanamide). \emph{-een}. \omterm{def:g11:functional-groups:carboxylic-acid}{Carbonzuren} (\ce{C=O} en
\ce{-OH} op hetzelfde koolstofatoom), \omterm{def:g11:functional-groups:carboxylic-acid}{esters} (\ce{C=O} en \ce{-O-}) en
\omterm{def:g11:functional-groups:amine}{amiden} (\ce{C=O} en \ce{N}): drie families.
\end{solution}

\begin{solution}{exo:g11:functional-groups:12}
Aspirine: een carboxylgroep (\omterm{def:g11:functional-groups:carboxylic-acid}{carbonzuur}) en een estergroep.
Paracetamol: een fenol-\ce{-OH} op de ring en een amidegroep
\ce{NH-CO}.
\end{solution}

\begin{solution}{exo:g11:functional-groups:13}
\ce{C2H6O}. Alleen ethanol, met zijn \ce{O-H}, vormt \omterm{def:g11:polarity-and-cohesion:hydrogen-bond}{waterstofbruggen}
tussen zijn \omterm{def:g7:atoms-and-molecules:molecule}{moleculen} (methoxy\-methaan heeft geen waterstof aan zijn
zuurstof). Ethanol kookt hoger, bij \qty{78}{\celsius}; methoxy\-methaan
bij \qty{-25}{\celsius}.
\end{solution}

\begin{solution}{exo:g11:functional-groups:14}
Een hydroxylgroep en een carboxylgroep. De keten heeft drie
koolstofatomen; het zuurkoolstofatoom is koolstofatoom 1 en de \ce{-OH}
zit op koolstofatoom 2: 2-hydroxypropaanzuur. Glycine:
2-aminoethaanzuur (vaak geschreven als aminoethaanzuur).
\end{solution}

\begin{solution}{exo:g11:functional-groups:15}
Ethanol \ce{CH3-CH2-OH}, \ce{C2H6O}, \omterm{def:g11:functional-groups:alcohol}{alcohol}; ethanal \ce{CH3-CHO},
\ce{C2H4O}, \omterm{def:g11:functional-groups:carbonyl}{aldehyde}; ethaanzuur \ce{CH3-COOH}, \ce{C2H4O2},
\omterm{def:g11:functional-groups:carboxylic-acid}{carbonzuur}. Van ethanol naar ethanal gaan twee waterstofatomen verloren;
van ethanal naar ethaanzuur komt er één zuurstofatoom bij.
\end{solution}

\begin{solution}{pb:g11:functional-groups:1}
\textbf{1.} A: estergroep; B: hydroxyl; C: carboxyl; D: carbonyl aan het
uiteinde van de keten; E: estergroep.

\textbf{2.} A \omterm{def:g11:functional-groups:carboxylic-acid}{ester}, B \omterm{def:g11:functional-groups:alcohol}{alcohol}, C \omterm{def:g11:functional-groups:carboxylic-acid}{carbonzuur}, D \omterm{def:g11:functional-groups:carbonyl}{aldehyde}, E \omterm{def:g11:functional-groups:carboxylic-acid}{ester}.

\textbf{3.} A (banaan) en E (ananas): fruitige geuren.

\textbf{4.} Primair: het koolstofatoom met de \ce{-OH} is gebonden aan
één koolstofatoom.

\textbf{5.} \ce{C6H12O}; bijvoorbeeld hexaan-2-on,
\ce{CH3-CO-CH2-CH2-CH2-CH3}.

\textbf{6.} De keten met de \ce{C-OH} heeft vier koolstofatomen,
genummerd vanaf de \ce{-OH}, met een methylgroep op koolstofatoom 3:
3-methylbutaan-1-ol.

\textbf{7.} C: ethaanzuur; D: hexanal.

\textbf{8.} Ethylbutanoaat, verwant aan butaanzuur en ethanol.

\textbf{9.} \emph{Ethanoaat}: het zuurdeel \ce{CH3-CO}, van ethaanzuur;
\emph{3-methylbutyl}: het alcoholdeel, van 3-methylbutaan-1-ol.

\textbf{10.} C (ethaanzuur) en B (3-methylbutaan-1-ol).

\textbf{11.} A \ce{C7H14O2}; B \ce{C5H12O}; C \ce{C2H4O2}.

\textbf{12.} $\ce{C2H4O2} + \ce{C5H12O}$ bevat \ce{C7H16O3}, één \ce{H2O}
meer dan \ce{C7H14O2}: water.

\textbf{13.} \ce{C2H4O2 + C5H12O -> C7H14O2 + H2O}: C 7 = 7, H 16 = 16,
O 3 = 3.

\textbf{14.} $M(\text{B}) = 5 \times 12.0 + 12 \times 1.0 + 16.0 =
\qty{88.0}{g/mol}$; $M(\text{C}) = \qty{60.0}{g/mol}$.

\textbf{15.} $60.0 + 88.0 - 18.0 = \qty{130.0}{g/mol}$.

\textbf{16.} $7 \times 12.0 + 14 \times 1.0 + 2 \times 16.0 =
\qty{130.0}{g/mol}$: hetzelfde.

\textbf{17.} E, \ce{C6H12O2}: $6 \times 12.0 + 12 \times 1.0 + 2 \times
16.0 = \qty{116.0}{g/mol}$.

\textbf{18.} \ce{C7H14O2}: dezelfde formule als A, dus een \omterm{def:g11:organic-skeletons:isomer}{isomeer}, met
dezelfde \omterm{def:g10:the-mole:molar-mass}{molaire massa}. Een \omterm{def:g10:the-mole:molar-mass}{molaire massa} alleen kan \omterm{def:g11:organic-skeletons:isomer}{isomeren} niet van
elkaar onderscheiden.

\textbf{19.} A en heptaanzuur hebben dezelfde \omterm{def:g7:atoms-and-molecules:atom}{atomen} maar verschillende
groepen, een \omterm{def:g11:functional-groups:carboxylic-acid}{ester} en een \omterm{def:g11:functional-groups:carboxylic-acid}{carbonzuur}: het ene ruikt naar banaan, het
andere ranzig. De groep, en waar die zit, beslist.

\textbf{20.} \qty{130.0}{g/mol}.
\end{solution}
